\subsection{Testing for Arbitrage in Generated Data}
\label{sec:arbitrage}

Static absence of arbitrage of each decoded surface does not by itself imply
absence of dynamic arbitrage between trading times. We here brute-force test a subset of possible arbitrage-strategies and will show that
while DYSANOS does not have arbitrage opportunities when traded options expire at their expiry -- as expected --, it does have arbitrage opportunities
if the same options can be sold at a later point for their mark to market.\\

\noindent
\textbf{Generators:}
We compare DYSANOS at smoothness~$\mu=0.5$ with a PCA-IV benchmark which
applies PCA directly to log-implied volatilities on the same regular grid  without any static arbitrage considerations. We therefore fully expect
it to have static arbitrage. We use this model as a canary -- any test for dynamic arbitrage should trigger in this model.\\


\noindent
\textbf{Scope:}
For each generator we use returns to expiry,
daily and weekly trading times, and
costs
corresponding to zero, one, and ten basis points. As spot is tradable to each option expiry, our experiments
include delta hedging.\\

\noindent
\textbf{Generated Market:}
We generate 100,000 and 1,000,000 samples for 10 trading times, daily or weekly. For the weekly generator, we also simulate
intermediate steps in order to obtain consistent spot values for options which expire before the next trading time. However, no
additional trading may occur at those times.

We use
the exact Gaussian transition for the autonomous surface state~$h$ and draw
spot under the importance-sampling law~$\O$ discussed in
Section~\ref{sec:importance}. If $p(\omega)$ and $q(\omega)$ are the densities
of a complete path under the fitted law~$\P$ and~$\O$, respectively, then
\begin{equation}\label{eq:arbitrage_importance_mixture}
 q(\omega)=0.1p(\omega)+0.9\bar q(\omega),
 \qquad
 D(\omega)=\frac{p(\omega)}{q(\omega)},
 \qquad
 \E_\P[F]=\E_\O[DF].
\end{equation}
The first mixture component draws the complete path directly from~$\P$; we
refer to these observations as ``direct-$\P$'' paths. The widened components use
a spot standard-deviation multiplier of eight and are tuned to place at least
1,000 observations on both sides of each relevant strike and return-sign
boundary. The reported runs satisfy $\E_\O[D]\simeq1$, and weighted spot and
return moments recover their direct-$\P$ estimates.\\

\noindent
\textbf{Trading Opportunities:}
We will consider trading opportunities conditional on each state: at each of our ten daily or weekly trading time $t$ conditional
on our state $x_t=(s_t,h_t)$, we are able to
trade calls for $M$ expiries and $N$
strikes per expiry, giving a total of $L=MN$ calls indexed by $\ell=1,\ldots,L$. 
In the
reported experiments, $M=7$, $N=20$, and $L=140$. Their surface coordinates are the
time-to-expiries and relative strikes $(T^\ell_t,K^\ell_t)$.
 The corresponding
fixed contracts have calendar expiries and cash strikes
\begin{equation}\label{eq:fixed_contract_coordinates}
 \tau^\ell_t=t+T^\ell_t,
 \qquad
 k^\ell_t=S_tK^\ell_t,
 \qquad
 C_t^\ell=S_tC(h_t;T^\ell_t,K^\ell_t),
\end{equation}
with Black\&Scholes delta~$\Delta_t^\ell$.
We consider holding either to the next trading time~$u$ or beyond expiry of each option, denoted by $u=\infty$. Spots between trading
times will be simulated as necessary to compute expiry values, no trading may occur at those times.


The return of each option from trading time $t$ to $u$ is then given as
\begin{equation}\label{eq:perdiexpriret}
 R_t^\ell
 := \left\{
 \begin{array}{ll}
	 S_u  C_u\left( u -t, \frac{ K^\ell_t }{ S_u } \right) -C_t^\ell & u\geq \tau^\ell_t \ ,\\
	 \left( S_{\tau^\ell_t} - S_t \right)^+  -C_t^\ell & u < \tau^\ell_t \ .
	\end{array}
	\right .
\end{equation}
In order to allow for delta hedging up to each expiry we add artificial strikes of zero to our list of strikes.
The zero strike call price return over the respective period is simply the spot return. As a result, we have $N':=N+1$ strikes
and $L':=N'M$ call contracts.

While we do not charge exit cost, we will charge proportional entry cost
\begin{equation}\label{eq:multi_expiry_costs}
 c_t^{\ell}=\gamma|C_t^\ell| \in \R^{L'}_{\geq 0}
\end{equation}
for cost rates~$\gamma\in\{0,10^{-4},10^{-3}\}$.


Let $a_t=a_t(s_t,h_t)\in\R^{L'}$ be a trading action. The P\&L from trading time $t$ to $u$ is given as
\begin{equation}\label{eq:multi_expiry_group_pnl}
 \pi_t(a_t)
 =
  \sum_{\ell=1}^{L'}
   a_t^\ell R_t^\ell-|a_t^\ell|c_t^\ell \ .
\end{equation}
We note that the zero strike calls introduced early allow the optimizer to construct a delta hedge,
or any other stock hedge, without defining a separate delta-hedged return explicitly.

A~portfolio~$a$ is an arbitrage opportunity for a given state $x_t = (s_t,h_t)$
if its P\&L is non-negative and positive for some samples.
It is clear that if $a$ is an arbitrage opportunity, then so is $a/|a|$. For numerical implementation we therefore impose~$\| a_t \|_1\leq 1$.


\subsubsection{Returns to Expiry}

In this section we test for each state $x_t$ on each sample path
whether their portfolios held to expiry ($u=\infty$) can generate arbitrage. Theoretically, this is not possible for DYSANOS as for each state all option
prices are free of static arbitrage; however, insufficient sampling can still generate such opportunities. An example is a very out of the money
put with a positive price, which simply does not pay on the generated set of paths. That means for DYSANOS the test in this section
validates the use of importance sampling for the spot.

We perform the following tests
\begin{enumerate}
	\item We first solve per state $x_t$ per sample $\omega$ the linear program of finding $a_t$ such that $\max \pi_t(a_t) > 0$ and $\min \pi_t(a_t) \geq 0$
	using synthetic spot prices (since $\pi$ is piece-wise affine this can be done on a much reduced synthetic spot set). 
	
	The resulting ``candidates'' $a_t(\omega)$ are real arbitrage opportunities in theory.

DYSANOS produces no candidates in this test, as expected.
	
	\item As a next step, we test whether these opportunities appear in sampled data.

We standardize $(s_t,h_t)$ at each trading
time using means and standard deviations computed from direct-$\P$ (non-importance sampled) paths, and
select the 1,024 nearest other states from either the original sample space (``Direct-$\P$'') or including the importance sampled space (``Weighted-$\P$'').
We compute the corresponding mean and standard errors and report in how many cases the mean exceeds three times its standard error.\footnote{
Details: with
$\bar D_r=D_r/\sum_vD_v$, the importance-weighted estimate and its
finite-sample standard error are
\begin{equation}\label{eq:multi_expiry_weighted_mean}
 \widehat\mu_D=\sum_r\bar D_r\pi_r,
 \qquad
 \operatorname{se}_D
 =\left\{
   \frac{\sum_r\bar D_r^2(\pi_r-\widehat\mu_D)^2}
        {1-\sum_r\bar D_r^2}
   \right\}^{1/2}.
\end{equation}
This estimate uses all neighbours drawn under~$\O$ and weights them by
$D=p/q$; we call it the \emph{importance-weighted $\P$ estimate}. Let
$\mathcal I_{\rm dir}$ be the neighbours belonging to the direct-$\P$
mixture component and $n_{\rm dir}:=|\mathcal I_{\rm dir}|$. The independent
\emph{direct-$\P$ estimate} is
\begin{equation}\label{eq:multi_expiry_direct_mean}
 \widehat\mu_{\rm dir}
 =\frac{1}{n_{\rm dir}}\sum_{r\in\mathcal I_{\rm dir}}\pi_r,
 \qquad
 \operatorname{se}_{\rm dir}
 =\left\{
   \frac{1}{n_{\rm dir}(n_{\rm dir}-1)}
   \sum_{r\in\mathcal I_{\rm dir}}(\pi_r-\widehat\mu_{\rm dir})^2
  \right\}^{1/2}.
\end{equation}
Thus both estimators target the same conditional mean under~$\P$: one uses
only paths drawn directly from~$\P$, while the other uses every neighbour
with its likelihood ratio. The tables report separately whether
$\widehat\mu_{\rm dir}>3\operatorname{se}_{\rm dir}$ and
$\widehat\mu_D>3\operatorname{se}_D$, as well as their intersection. These
conditional-mean calculations measure the physical incidence and scale of a
complete-support certificate; they do not replace its pathwise non-loss
condition.

For each universe we select ten disjoint batches of 1,000 paths. Their ten
trading times give 100,000 separately solved state--time LPs in every row
of Tables~\ref{tab:multi_expiry_100k} and~\ref{tab:multi_expiry_1m}. Increasing
the universe from 100,000 to one million paths therefore changes the sampled
states and the population available to the nearest-neighbour diagnostic, but
not the table denominator.}
\end{enumerate}






\begin{table}[H]
\centering
\scriptsize
\setlength{\tabcolsep}{2pt}
\begin{tabular}{llrrrrrrrr}
\toprule
 & & \multicolumn{4}{c}{PCA-IV occurrence (\%)}
       & \multicolumn{4}{c}{DYSANOS occurrence (\%)}\\
\cmidrule(lr){3-6}\cmidrule(lr){7-10}
Trading interval & Cost & Candidates & Direct-$\P$ & Weighted-$\P$ & Both
                         & Candidates & Direct-$\P$ & Weighted-$\P$ & Both\\
\midrule
Daily  & $0$ bp  & $5.273\%$ & $2.333\%$ & $3.282\%$ & $1.721\%$ & $0.000\%$ & $0.000\%$ & $0.000\%$ & $0.000\%$\\
Daily  & $1$ bp  & $0.950\%$ & $0.398\%$ & $0.743\%$ & $0.331\%$ & $0.000\%$ & $0.000\%$ & $0.000\%$ & $0.000\%$\\
Daily  & $10$ bp & $0.042\%$ & $0.037\%$ & $0.042\%$ & $0.037\%$ & $0.000\%$ & $0.000\%$ & $0.000\%$ & $0.000\%$\\
Weekly & $0$ bp  & $5.525\%$ & $2.520\%$ & $3.712\%$ & $2.043\%$ & $0.000\%$ & $0.000\%$ & $0.000\%$ & $0.000\%$\\
Weekly & $1$ bp  & $0.923\%$ & $0.466\%$ & $0.750\%$ & $0.415\%$ & $0.000\%$ & $0.000\%$ & $0.000\%$ & $0.000\%$\\
Weekly & $10$ bp & $0.059\%$ & $0.057\%$ & $0.059\%$ & $0.057\%$ & $0.000\%$ & $0.000\%$ & $0.000\%$ & $0.000\%$\\
\bottomrule
\end{tabular}
\caption{Statewise multi-expiry arbitrage test with a 100,000-path universe.
Every percentage uses 100,000 state--time pairs as its denominator.
``Candidates'' is the occurrence of theoretical arbitrage opportunities. ``Direct-$\P$'' and
``Weighted-$\P$'' require, respectively, a conditional mean exceeding three
standard errors under the direct and importance-weighted estimators; ``Both''
is their intersection.}
\label{tab:multi_expiry_100k}
\end{table}

\begin{table}[H]
\centering
\scriptsize
\setlength{\tabcolsep}{2pt}
\begin{tabular}{llrrrrrrrr}
\toprule
 & & \multicolumn{4}{c}{PCA-IV occurrence (\%)}
       & \multicolumn{4}{c}{DYSANOS occurrence (\%)}\\
\cmidrule(lr){3-6}\cmidrule(lr){7-10}
Trading interval & Cost & Candidates & Direct-$\P$ & Weighted-$\P$ & Both
                         & Candidates & Direct-$\P$ & Weighted-$\P$ & Both\\
\midrule
Daily  & $0$ bp  & $5.097\%$ & $2.158\%$ & $3.022\%$ & $1.564\%$ & $0.000\%$ & $0.000\%$ & $0.000\%$ & $0.000\%$\\
Daily  & $1$ bp  & $0.844\%$ & $0.303\%$ & $0.618\%$ & $0.258\%$ & $0.000\%$ & $0.000\%$ & $0.000\%$ & $0.000\%$\\
Daily  & $10$ bp & $0.043\%$ & $0.029\%$ & $0.043\%$ & $0.029\%$ & $0.000\%$ & $0.000\%$ & $0.000\%$ & $0.000\%$\\
Weekly & $0$ bp  & $5.537\%$ & $2.510\%$ & $3.689\%$ & $1.970\%$ & $0.000\%$ & $0.000\%$ & $0.000\%$ & $0.000\%$\\
Weekly & $1$ bp  & $0.940\%$ & $0.426\%$ & $0.767\%$ & $0.375\%$ & $0.000\%$ & $0.000\%$ & $0.000\%$ & $0.000\%$\\
Weekly & $10$ bp & $0.059\%$ & $0.050\%$ & $0.059\%$ & $0.050\%$ & $0.000\%$ & $0.000\%$ & $0.000\%$ & $0.000\%$\\
\bottomrule
\end{tabular}
\caption{Statewise multi-expiry arbitrage test with a one-million-path
universe. Every percentage uses the same 100,000 audited state--time
denominator and definitions as Table~\ref{tab:multi_expiry_100k}.}
\label{tab:multi_expiry_1m}
\end{table}

We run this experiment for 100,000 and 1,000,000 paths.
The larger universe leaves the conclusion unchanged. PCA-IV certificate
occurrence rates are stable across path counts. The occurrence significant
under both physical estimators is somewhat lower in
the larger universe, but remains positive in every case. We retain both
tables to display this difference. PCA-IV contains complete-support arbitrage
at zero cost and after one- and ten-basis-point costs.

DYSANOS produces no certificate in any reported to-expiry case. The PCA-IV detections are
the positive control: the same statewise test finds its known payoff
violations without imposing a particular monotonicity, convexity, or calendar
spread in advance. Within the sampled states and the tested to-expiry strategy
class, the results therefore detect conditional arbitrage in PCA-IV and do
not detect it in DYSANOS.

\subsubsection{Period Returns}

The preceding test holds each option to expiry. To test trading between
let~$u$ be the next trading time in~\eqref{eq:perdiexpriret}, conditional
on neighbouring states.


We split our paths, stratified by importance-sampling component, into
$80\%$ training and $20\%$ validation sets. For each~$t$, the state
coordinates are standardized using only the training set. We select 1,024
training states $(\bar s,\bar h)$ from the direct-$\P$ component. Let
$\mathcal N_t(\bar s,\bar h)$ contain the nearest training states in the
standardized Euclidean metric: 1,024 neighbours for 100,000 paths and 4,096
for one million paths. At each selected state we solve
\begin{equation}\label{eq:period_local_lp}
\begin{split}
 \max_{a_t}\quad
 &\sum_{\omega\in\mathcal N_t(\bar s,\bar h)}
   \bar D_\omega\pi_t^\omega(a_t),\\
 \text{subject to}\quad
 &\|a_t\|_1\leq1,\\
 &\pi_t^\omega(a_t)\geq0,
 \qquad \omega\in\mathcal N_t(\bar s,\bar h),
\end{split}
\end{equation}
where the importance weights~$\bar D_\omega$ are normalized within the
training neighbourhood. This gives a locally constant approximation to the
conditional action $a_t=a_t(s_t,h_t)$ without restricting it to individual
options or prescribed spreads.

The action, state normalization, center, and radius of the smallest closed
ball containing its training neighbours are then frozen. Validation uses all
paths in the disjoint validation set whose states lie in that ball. Let $\eta(a_t)$
be a numerical threshold.

A validation observation is a loss when
$\pi_t(a_t)<-\eta_t(a_t)$ for each $\omega$ in the validation ball, and a gain when
$\pi_t(a_t)>\eta_t(a_t)$. A portfolio is considered an arbitrage only when there is
no loss in its validation ball and the lower confidence bounds for both mean
P\&L and gain probability are positive under both the direct-$\P$ sample and
the importance-weighted sample. We use the Bonferroni level
$0.05/442{,}368$ for these simultaneous statements. A candidate is ``unresolved''
if its validation ball contains fewer than twenty direct-$\P$ paths or has
importance-weighted effective sample size below twenty. 

There are nine period returns and 1,024
selected states per return, so every occurrence rate below has denominator
$9{,}216$. Unlike the complete-support test above, this finite-neighbourhood
test is an empirical detection of conditional arbitrage within the stated
state regions.

\begin{table}[H]
\centering
\small
\setlength{\tabcolsep}{4pt}
\begin{tabular}{lrrrrr}
\toprule
 & & \multicolumn{2}{c}{PCA-IV occurrence (\%)}
       & \multicolumn{2}{c}{DYSANOS occurrence (\%)}\\
\cmidrule(lr){3-4}\cmidrule(lr){5-6}
Trading interval & Cost & Arbitrage & Unresolved
                         & Arbitrage & Unresolved\\
\midrule
Daily  & $0$ bp  & $0.803\%$ & $10.569\%$ & $7.107\%$ & $9.429\%$\\
Daily  & $1$ bp  & $0.315\%$ & $10.569\%$ & $5.534\%$ & $9.429\%$\\
Daily  & $10$ bp & $0.054\%$ & $1.421\%$  & $2.365\%$ & $9.429\%$\\
Weekly & $0$ bp  & $0.130\%$ & $7.975\%$  & $2.268\%$ & $7.227\%$\\
Weekly & $1$ bp  & $0.076\%$ & $7.975\%$  & $1.736\%$ & $7.227\%$\\
Weekly & $10$ bp & $0.022\%$ & $5.968\%$  & $0.738\%$ & $7.140\%$\\
\bottomrule
\end{tabular}
\caption{Fixed-contract period-return test with 100,000 paths.
 ``Unresolved'' occurrences lack the prescribed validation support and are
not classified as failures.}
\label{tab:period_comparison_100k}
\end{table}

The 100,000-path test retains finite-neighbourhood candidates in every case.
DYSANOS has the higher validated occurrence in every row. Higher transaction
costs reduce occurrence. Both generators retain substantial unresolved
regions, so the larger universe is needed to distinguish a stable finding
from sparse validation support.

\begin{table}[H]
\centering
\small
\setlength{\tabcolsep}{4pt}
\begin{tabular}{lrrrrr}
\toprule
 & & \multicolumn{2}{c}{PCA-IV occurrence (\%)}
       & \multicolumn{2}{c}{DYSANOS occurrence (\%)}\\
\cmidrule(lr){3-4}\cmidrule(lr){5-6}
Trading interval & Cost & Arbitrage & Unresolved
                         & Arbitrage & Unresolved\\
\midrule
Daily  & $0$ bp  & $1.584\%$ & $0.000\%$ & $12.500\%$ & $0.000\%$\\
Daily  & $1$ bp  & $0.564\%$ & $0.000\%$ & $9.082\%$  & $0.000\%$\\
Daily  & $10$ bp & $0.141\%$ & $0.000\%$ & $3.266\%$  & $0.000\%$\\
Weekly & $0$ bp  & $0.423\%$ & $0.000\%$ & $2.365\%$  & $0.000\%$\\
Weekly & $1$ bp  & $0.358\%$ & $0.000\%$ & $1.758\%$  & $0.000\%$\\
Weekly & $10$ bp & $0.184\%$ & $0.000\%$ & $0.412\%$  & $0.000\%$\\
\bottomrule
\end{tabular}
\caption{Fixed-contract period-return test with one million paths. The
definitions, denominator, and cost hierarchy are those of
Table~\ref{tab:period_comparison_100k}.}
\label{tab:period_comparison_1m}
\end{table}

The larger universe resolves every selected state--time center. The retained
occurrence changes materially in several rows. DYSANOS has the higher
occurrence in every row, and both generators retain candidates at ten basis
points. Thus the period-return detections are not explained by sparse
validation support in the smaller universe.

The neighbourhood test can miss a rare loss. We therefore perform a second
numerical test on the two strongest DYSANOS weekly raw candidates at ten basis
points, ranked by the smaller of the direct-$\P$ and importance-weighted lower
mean bounds. For each candidate, we freeze its exact initial state and
unit-$\ell^1$ action, then draw one million independent five-day
continuations under the fitted physical law~$\P$. Contracts expiring within
the week and their matched spot coordinates close at the exact expiry; all
other positions close after five daily steps. If $\pi_r$ and $\eta_r$ are the
resulting P\&L and numerical band, we record
\begin{equation}\label{eq:period_exact_state_loss}
 \widehat p_-=\frac{1}{N}\sum_{r=1}^N\bm 1_{\{\pi_r<0\}},
 \qquad
 \widehat p_{-,\eta}=\frac{1}{N}\sum_{r=1}^N
                 \bm 1_{\{\pi_r<-\eta_r\}},
 \qquad N=10^6.
\end{equation}
The observed loss occurrences and P\&Ls are reported in
Table~\ref{tab:period_exact_state_replay}.

\begin{table}[H]
\centering
\small
\setlength{\tabcolsep}{5pt}
\begin{tabular}{rrrrrr}
\toprule
Rank & Loss (\%) & Beyond band (\%) & Minimum P\&L
     & Mean P\&L & MC error\\
\midrule
$1$ & $0.0000\%$ & $0.0000\%$ & $ 2.963\times10^{-4}$
    & $4.683\times10^{-3}$ & $1.019\times10^{-6}$\\
$2$ & $0.0752\%$ & $0.0746\%$ & $-1.478\times10^{-3}$
    & $5.821\times10^{-3}$ & $2.175\times10^{-6}$\\
\bottomrule
\end{tabular}
\caption{Exact-state conditional resampling of the two strongest DYSANOS
weekly raw period-return candidates at ten basis points. Each row uses one
million fresh physical continuations of the frozen state and action.}
\label{tab:period_exact_state_replay}
\end{table}

The first portfolio uses one expiry-matched stock leg and eight calls. It is
long the two-day stock leg, long calls at ten and twenty days, and short calls
at two, five, forty days, and six months. The second is call-only: it combines
small front-end call spreads and long forty-day and six-month calls with a
dominant short one-year in-the-money call and a small long one-year
at-the-money call. The
second candidate is rejected by observed losses beyond the numerical band.
The first has no loss in one million continuations; its one-sided $95\%$
Clopper--Pearson upper bound on loss probability is
$1-0.05^{1/N}=2.996\times10^{-6}$, or $0.000300\%$. It is not falsified by
this numerical test, but finite sampling does not certify non-negative P\&L
on the complete conditional support.

The zero to-expiry counts above follow from the cross-sectional SANOS
constraints on each decoded surface and its terminal payoff cone. Those
constraints do not restrict the transition
$C_u(\tau,k)-C_t(\tau,k)$ of the same fixed contract. The period-return results
therefore identify state regions requiring a separate conditional test. The
exact-state replay rejects one of the two strongest DYSANOS candidates and
leaves the other unresolved at the stated sampling precision. It does not
establish non-negative P\&L on the complete conditional support, but, considering
the evidence, is a strong indicator that DYSANOS does admit dynamic arbitrage.


\endinput
